\documentclass[output=paper,colorlinks,citecolor=brown]{langscibook}
\ChapterDOI{10.5281/zenodo.20611905}
\author{Viktor Köhlich\orcid{}\affiliation{Goethe University Frankfurt}}

\title{The uninflecting word class \textit{rentaishi} in Modern Japanese}

\abstract{In this paper, I focus on the Japanese word class \emph{rentaishi} (`noun-modifying words'), which is often overlooked in Western linguistics but has typological significance in the debate on uninflectedness. While Japanese verbs and adjectives exhibit rich inflectional morphology for categories like tense, aspect, and modality, \emph{rentaishi} are restricted to attributive position, and lack inflection completely. Furthermore, they display morphological diversity, contrary to the major Japanese word classes (e.g., \emph{aru} ‘a certain’, \emph{rei-no} ‘mere’). Historically, the word class \emph{rentaishi} emerged in the early 20\textsuperscript{th} century when Japanese linguists recognized that Japanese adjectives inflect for tense and can appear predicatively independently, unlike adjectives in Western languages that require a copula, both to express tense and to appear predicatively, and must agree with the head noun in case, number and gender in attributive position. To account for non-predicative, attributive modifiers that do not inflect, Japanese linguists introduced the word class \emph{rentaishi}, grouping syntactically restricted and morphologically deficient lexemes. As a future consideration, I explore whether \emph{rentaishi} should remain a distinct word class in Japanese or if relevant lexemes should be reclassified as defective members of other word classes or even as part of the adjective group they were originally separated from, presenting arguments for both approaches.}


\IfFileExists{../localcommands.tex}{
   \addbibresource{../localbibliography.bib}
   \usepackage{tabularx,multicol}
\usepackage[hyphens]{url}
\urlstyle{same}

\usepackage{langsci-optional}
\usepackage{langsci-lgr}
\usepackage{langsci-gb4e}

\usepackage{soul}
% \usepackage{booktabs}
% \usepackage{geometry}
\usepackage{multirow}
% \usepackage{makecell} %safely breaks a line inside a table cell
% \usepackage{tablefootnote} %footnotes in table captions
\usepackage{enumerate}
\usepackage{tikz}
\usepackage{array}
\usepackage[shortlabels]{enumitem}

%for having several figures on one line, turning words in 90°angles
% \usepackage{graphicx}
\usepackage{subcaption}

% \usepackage{caption} %for making the caption line longer in concrete cases

% \usepackage{rotating} %rotate a whole page; for large tables

\usepackage{fontspec}
\newfontfamily\charis{Charis SIL} %for being able to use the symbol ‿ and a better command of some of the boldfaced diacritics in 01
\newfontfamily\japanesefont{Noto Serif CJK JP} % for Japanese script
% \usepackage{tipa}

% \usepackage{needspace} %to keep exs on one page

\usepackage{xcolor}

\usepackage{pgfplots}
\pgfplotsset{compat=1.18}

%%%%%   AVMs for 02	%%%%%
\usepackage{langsci-avm}
\avmsetup{switch = ?}


\usepackage{langsci-branding}

   %\newcommand*{\orcid}{}


\makeatletter
\let\thetitle\@title
\let\theauthor\@author
\makeatother

\newcommand{\togglepaper}[1][0]{
   \bibliography{../localbibliography}
   \papernote{\scriptsize\normalfont
     \theauthor.
     \titleTemp.
     To appear in:
     E. Di Tor \& Herr Rausgeberin (ed.).
     Booktitle in localcommands.tex.
     Berlin: Language Science Press. [preliminary page numbering]
   }
   \pagenumbering{roman}
   \setcounter{chapter}{#1}
   \addtocounter{chapter}{-1}
}

\newbool{bookcompile}
\booltrue{bookcompile}
\newcommand{\bookorchapter}[2]{\ifbool{bookcompile}{#1}{#2}}



\renewcommand{\thempfootnote}{\fnsymbol{mpfootnote}} % use * for footnotes in float env

% \newcolumntype{L}[1]{>{\raggedright\arraybackslash}p{#1}} %left-aligned (non-justified) text in tables
% \newcolumntype{R}[1]{>{\raggedleft\arraybackslash}p{#1}} % right-aligned (non-justified) text in tables

\definecolor{customgreen}{RGB}{27,158,119}
\definecolor{custompink}{RGB}{231,42,138}
\definecolor{customblue}{RGB}{117,112,179}
\definecolor{customlightblue}{RGB}{143,170,220}
\definecolor{customdarkblue}{RGB}{68,114,196}
\definecolor{customyellow}{RGB}{225,192,0}
\definecolor{customorange}{RGB}{233,137,21}
\definecolor{customgray}{RGB}{229,229,229}


% Left-aligned indexes
% % \makeatletter
% % \patchcmd{\theindex}{\twocolumn}{\raggedright\twocolumn}{}{}
% % \makeatother

%to make Japanese scripts in the bibliography smaller
\newcommand{\japbib}[1]{{\japanesefont\small #1}}

%%%%%%%%%%%%%%%%%%%	MACROS for 02	%%%%%%%%%%%%%%%%

%%% Formatting	%%%%%

\newcommand{\textex}[1]{\textit{#1}} 		 %for formatting linguistic examples in-text%
%command-shift X

\newcommand{\lxm}[1]{\textsc{#1}}  		%for formatting names of lexemes
%command-option-shift L

%%%%%%%%% Abbreviations	%%%%%%

\newcommand{\featname}[1]{\textsc{#1}}
\newcommand{\featspec}[2]{\featname{#1}:{#2}}
\newcommand{\feat}[2]{\textsc{#1}:{#2}}  				% \feat{FEATURE}{value}  sc's
\newcommand{\featnameonly}[1]{\textsc{#1}} 			% \featnameonly{featurename}  sc's

\newcommand{\pr}{$^{\prime}$}

\newcommand{\Corr}{\textit{Corr}}
\newcommand{\propm}{\textit{pm}}

\newcommand{\PC}{Π$_{C}$}

\newcommand{\PF}{Π$_{F}$}
\newcommand{\PR}{Π$_{R}$}

\newcommand{\lex}{$\mathcal{L}$}

   %% hyphenation points for line breaks
%% Normally, automatic hyphenation in LaTeX is very good
%% If a word is mis-hyphenated, add it to this file
%%
%% add information to TeX file before \begin{document} with:
%% %% hyphenation points for line breaks
%% Normally, automatic hyphenation in LaTeX is very good
%% If a word is mis-hyphenated, add it to this file
%%
%% add information to TeX file before \begin{document} with:
%% %% hyphenation points for line breaks
%% Normally, automatic hyphenation in LaTeX is very good
%% If a word is mis-hyphenated, add it to this file
%%
%% add information to TeX file before \begin{document} with:
%% \include{localhyphenation}
\hyphenation{
    Al-ex-an-der
    Ber-ber
    chan-ges
    cir-cum-stan-ces
    coin-age
    Fraj-zyngier
    ita-liano
    Ja-pa-nese
    Kö-nig
    Krzy-żanowski
    li-te-ra-tur-no-go
    mi-zen-kei
    par-a-digm
    phe-nom-e-non
    ren-tai-shi
    Slo-vene
    Spen-cer
    Stoj-kov
    Ver-bi-ca-re-se
    wide-spread
}

\hyphenation{
    Al-ex-an-der
    Ber-ber
    chan-ges
    cir-cum-stan-ces
    coin-age
    Fraj-zyngier
    ita-liano
    Ja-pa-nese
    Kö-nig
    Krzy-żanowski
    li-te-ra-tur-no-go
    mi-zen-kei
    par-a-digm
    phe-nom-e-non
    ren-tai-shi
    Slo-vene
    Spen-cer
    Stoj-kov
    Ver-bi-ca-re-se
    wide-spread
}

\hyphenation{
    Al-ex-an-der
    Ber-ber
    chan-ges
    cir-cum-stan-ces
    coin-age
    Fraj-zyngier
    ita-liano
    Ja-pa-nese
    Kö-nig
    Krzy-żanowski
    li-te-ra-tur-no-go
    mi-zen-kei
    par-a-digm
    phe-nom-e-non
    ren-tai-shi
    Slo-vene
    Spen-cer
    Stoj-kov
    Ver-bi-ca-re-se
    wide-spread
}

   \boolfalse{bookcompile}
   \togglepaper[23]%%chapternumber
}{}

\begin{document}
\maketitle

\is{rentaishi|(}
\section{Introduction} 
In this paper, I discuss the not universally accepted word class of so-called \emph{rentaishi} in Modern Japanese. \emph{Rentaishi} are a group of lexemes that are syntactically restricted to their use as attributive modifiers. From a morphological perspective, this syntactic restriction translates to absence of inflection. Different instances of lexemes that have been called \emph{rentaishi} are given in (\ref{12ex1}) below, as part of a DP and highlighted in boldface (cf. \citealt{sekine1937, kai1980, kim2006, matsubara2009})\il{Japanese|(}.

\begin{exe}
	\ex \label{12ex1}
	\begin{xlist}
		\ex 
		\gll \textbf{aru} gakusei \\
		certain student \\
		\trans ‘a certain student’
		
		\ex 
		\gll \textbf{taishita} sa \\
		considerable difference \\
		\trans ‘a considerable difference’
		
		\ex 
		\gll \textbf{shutaru} gen’in \\
		main reason \\
		\trans ‘the main reason’
		
		\ex 
		\gll \textbf{ōkina} ie \\
		big house \\
		\trans ‘big house’
	\end{xlist}
\end{exe}

As visible above, every relevant lexeme serves as a modifier of a noun, and that is the only syntactic environment such lexemes appear in. Importantly, they all exhibit different morphological endings. This morphological heterogeneity differentiates them from the major inflecting word classes in Japanese, namely verbs, which end with the vowel \emph{-u}, and adjectives, which either consistently end with the element \emph{-i} (\emph{i}-adjectives) or with the element \emph{-na} in attributive and \emph{-da} in predicative position (\emph{na}-adjectives). This is due to the fact that many lexemes that have been classified as \emph{rentaishi} are derived from other word classes and have become frozen or fossilized.

In this paper, \emph{rentaishi} will be discussed from a perspective of inflection, or inflectability. Concretely, I will show that, historically, the emergence of the word class \emph{rentaishi} is tied to the criterion of \isi{uninflectedness} and, in fact, a direct result of the greater awareness exhibited by Japanese linguists at the beginning of the 20\textsuperscript{th} century concerning the differences between Japanese and the languages of the West, especially \ili{English} and \ili{Dutch}. The need to fill the slot of dependent attributive modifiers in Japanese as equivalents to the Western adjectives, which they perceived as defective contrary to the independent and inflecting Japanese adjectives (\emph{keiyōshi}), lead to the creation of the word class \emph{rentaishi}. This creation marks a linguistic turn at the beginning of the 20\textsuperscript{th} century.

However, there is more to be discussed about the role of these modifiers in the Japanese word class system from a modern perspective and in a cross-linguistic context in order to evaluate whether they deserve a status as an independent word class. For example, \citeauthor{muraki1996} (\citeyear{muraki1996} et seq.) has put forth a new analysis of \emph{rentaishi} as non-predicative adjectives, which marks relevant lexemes as \emph{uninflectable} rather than \emph{uninflecting} in the terms of \citet{spencer2020}. Towards the end of this paper, I will explore a possible re-analysis of relevant lexemes as uninflectable and defective, rather than uninflecting.

This paper is structured as follows. First, the relevant terminology will be discussed in Section \ref{12sec2}. In Section \ref{12sec3}, the notion of \emph{inflection} will be discussed from the perspective of Japanese linguistics, and the category of \emph{rentaishi} will be introduced. The main analysis is carried out in Section \ref{12sec4} with special reference to the historical dimension of this group and its relevance for the debate on \emph{\isi{uninflectedness}}. Section \ref{12sec5} then discusses this group from a modern perspective and Section \ref{12sec6} concludes the paper.\footnote{For all transliteration, I adopt the Hepburn-Romanization System for Japanese, even when citing from other sources. Japanese names are given with the last name first written in capital letters. I am very grateful to Ken Hiraiwa for judgments on the Modern Japanese data as well as to two anonymous editors for very useful comments and especially for pointing my attention towards the recently published PhD thesis of \citet{nespoli2023}. All remaining errors are, of course, my own.}

\section{Terminology} \label{12sec2}

\textit{Uninflectedness}\is{uninflectedness} refers to the inability of a given lexeme or a group of lexemes to inflect. A prerequisite for adopting this term is that the language under investigation, in principle, allows for inflection and that inflection has a syntactic dimension. \citet[33]{baermanbrowncorbett2005} define \isi{uninflectedness} as follows:

\begin{exe}
	\ex \label{12ex2}
	\begin{xlist}
		\exi{(i)} There is, in certain lexemes only, a loss of all values of a particular feature F found elsewhere in the language. This loss may depend on the presence of a particular combination of values of one or more other features (the context).
		\exi{(ii)} Other syntactic objects distinguish values of feature F, either generally or in the given context, and feature F is therefore syntactically relevant.
	\end{xlist}
\end{exe}

What this quote emphasizes is that \isi{uninflectedness} is a morphological phenomenon, the syntax is autonomous, as will become clearer below.

Furthermore, the notion \emph{\isi{uninflectedness}}, in the first place, denotes the general inability of a given lexeme to inflect, or a general lack of inflection, but in fact, it is multidimensional. A first important distinction can be drawn between \textit{uninflecting} and \textit{uninflectable} (groups of) lexemes. In this paper, I will follow \citet[142]{spencer2020} and \citet{baermanbrowncorbett2005, baermanbrowncorbett2017} and with the terminology \textit{uninflecting} refer to parts of speech and elements belonging to these parts of speech that never, under no circumstances, inflect in a given language. Depending on the language, those can be lexical or functional parts of speech. A typical example of an uninflecting word class\is{uninflecting!word class} are conjunctions as those very rarely inflect across the world's languages. Examples are \ili{Russian} \emph{li} `or', \ili{German} \emph{und} `and' or the equivalents in \ili{English}. On the other hand, the lexical class of nouns shows inflection in some languages, in this case makes up an \textit{inflecting} word class, but not in others, where nouns then are uninflecting.

On the contrary, the notion \textit{uninflectable}\is{uninflectable!element} refers to specific elements that belong to an otherwise inflecting, lexical or functional, part of speech in a given language, but which for some reason or other do not inflect, that is do not show the relevant endings and inflectional paradigms of other elements belonging to the same part of speech in the relevant syntactic environments. Crucially, uninflectable lexemes belong to an inflecting word class, hence are expected to inflect \citep[142]{spencer2020}. An important prerequisite for analyzing a lexeme as \emph{uninflectable} is therefore that it appears in all the relevant syntactic environments other lexemes from the same part of speech appear in \citep[142--144]{spencer2020}.

An often-quoted example of an uninflectable lexeme is the \ili{Russian} noun \emph{pol'to} `coat'. Nouns in \ili{Russian} inflect for case and number, therefore compose an inflecting word class. However, the noun \emph{pol'to} simply does not show any specific inflectional ending(s) for those categories, all forms being syncretic. In other words,  this noun displays the same inflectional form across all cells in its inflectional paradigm \citep[30]{baermanbrowncorbett2005}. Nevertheless, importantly, \emph{pol'to} can appear in all relevant syntactic environments in which nouns in \ili{Russian} typically appear, for example as the subject or object of a sentence. Its uninflectable nature becomes visible when a controller appears, such as the preposition \emph{v} `in'. As visible in (\ref{12ex3b}) below, when this preposition takes a DP containing a noun and an adjective as its complement, the adjective shows obligatory concord within the PP and is inflected for prepositional case. However, the noun \emph{pol'to} is unaltered and appears in the same morphological form as in nominative contexts, exemplified in (\ref{12ex3a}).

\begin{exe}
	\ex \label{12ex3}
	\begin{xlist}
		\ex \label{12ex3a}
		\gll bel-oe pol’to \\
		white-\textsc{nom.sg.n} coat.Ø \\
		\trans `white coat’
		
		\ex \label{12ex3b}
		\gll v bel-om pol’to \\
		in white-\textsc{prep.sg.n} coat.Ø \\
		\trans `in a white coat’
	\end{xlist}
\end{exe}

Lexemes can also be \textit{partially uninflectable}\is{uninflectable!element!partially (\emph{museum} nouns)}. This means they belong to an inflecting part of speech, but show a restricted inventory of inflectional forms across their paradigm. In other words, they show inflection in some contexts, but not in others. Staying in the nominal domain, \emph{context} refers to case, gender and number. For example, the \ili{Polish} noun \emph{muzeum} `museum' only shows distinct inflectional forms for case in its plural cells, but in singular all entries are syncretic (\citealt[30]{baermanbrowncorbett2005}; \citealt[143]{spencer2020}).

Finally, \textit{contextual uninflectability}\is{uninflectability!contextual} means that certain lexical items do not inflect in specific syntactic environments, but otherwise do \citep[149--150]{spencer2020}. This is for example the case for \ili{German} (and \ili{Dutch}) adjectives, which do inflect, but never in predicative position as shown below (cf. \citealt{roehrs2008, zeijlstra2020} among others).

\begin{exe}
	\ex \label{12ex4}
	\begin{xlist}
		\ex 
		\gll das rot-e Buch \\
		the red-\textsc{nom.sg.n} book.\textsc{nom.sg.n} \\
		\trans `the red book’
		
		\ex 
		\gll Das Buch ist rot. \\
		the book.\textsc{nom.sg.n} is red.\textsc{Ø} \\
		\trans `The book is red.’
	\end{xlist}
\end{exe}

Importantly, however, all the lexemes depicted above do appear in the relevant syntactic environments expected from their part of speech. If a lexeme fails to appear in the syntactic environments it is supposed to appear in in the first place, it can be called syntactically \textit{defective}\is{defectiveness!syntactical} \citep[143]{spencer2020}. In this sense, this notion can refer to adjectives which cannot appear predicatively for syntactic reasons. This is the case for so-called \emph{non-intersective} or \emph{intensional} adjectives\is{adjective!non-intersective} such as \emph{former} or \emph{alleged} and their cross-linguistic counterparts \citep{bolinger1967, kamppartee1995, larson1998}, or so-called relational adjectives\is{adjective!relational} (\citealt{levi1978, fabregas2007, bortolotto2016} a.o.). For instance, although adjectives in \ili{Russian}, different to \ili{German}, inflect in attributive as well as predicative position, the group of relational adjectives\is{adjective!relational}, such as \emph{kni\v{z}nyj} `pertaining to books' derived from \emph{kniga} `book', is simply illicit in predicative position as shown in (\ref{12ex5}) \citep{mezhevich2002, bailyn2002}.

\begin{exe}
	\ex \label{12ex5}
	\begin{xlist}
		\ex 
		\gll kni\v{z}-n-yj magazin \\
		book\textsc{-ra-nom.sg.m} store.\textsc{nom.sg.m} \\
		\trans ‘bookstore’ (lit. ‘store pertaining to books’)
		
		\ex[*] {
			\gll Magazin byl kni\v{z}-n-yj. \\
			store.\textsc{nom.sg.m} was book\textsc{-ra-nom.sg.m} \\
			\trans Intend.: ‘The store was booky/bookish.’ \hfill (\citealt[100]{mezhevich2002})}
	\end{xlist}
\end{exe}

The distinction between morphologically and syntactically restricted modifiers will become relevant for our discussion of \emph{rentaishi} and Japanese adjectives in Section \ref{12sec5} of this paper.

\section{Japanese} \label{12sec3}

This section gives an overview over what inflection denotes from a Japanese perspective, illustrating with different parts of speech. After this, the category of \emph{rentaishi} will be outlined with reference to its relevance to the present discussion.

\subsection{Overview} \label{12sec3.1}

Japanese is a strictly head-final SOV language in which the verb cannot be moved from sentence-final position (\citealt{kuno1973, martin1975, iwasaki2013} a.o.). Another key feature is its agglutinative character which means that grammatical information is encoded via affixes, predominantly suffixes. Concerning word classes, since the Edo-period (1603--1868), a classical divide has been drawn between \emph{inflecting} and \emph{uninflecting}, as well as \emph{dependent} and \emph{independent} word classes. \emph{Inflecting} word classes can be further differentiated by their distinct morphology.

\subsection{Inflection from a Japanese perspective} \label{12sec3.2}

While established linguistic notions such as \emph{inflection} are often based on European languages -- usually Latin or Ancient Greek -- the preceding outline has shown that \emph{inflection} has been highly relevant for differentiating word classes in Japanese. In fact, a similar notion was established by linguists in Japan during the Edo-period, namely the term \emph{katsuyō}\is{katsuyō|(} {\japanesefont 活用}. The linguistic notion of \emph{katsuyō} presumably goes back to MOTOORI Norinaga (1730--1801)\ia{Motoori, Norinaga}, one of the most important national linguists (\emph{kokugakusha}) of the Edo-period, and was included as a criterion to categorize Japanese parts of speech in the \emph{school grammar}\is{Japanese school grammar} by HASHIMOTO Shinkichi (1882--1945) (Hashimoto 1948 [1934])\ia{Hashimoto, Shinkichi}, so-called because it is to date the grammar taught in Japanese school curricula.

\subsubsection{Inflection of verbs} \label{12sec3.2.1}

It was the goal of the \emph{kokugakusha} to analyze the Japanese language independently of other languages. And while \emph{katsuyō} is in \ili{English} often translated as \emph{inflection}, it does not mean exactly the same thing. But then what is \emph{katsuyō}? As mentioned above, an important difference is drawn between inflecting and uninflecting parts of speech in Japanese. The Japanese terms are \emph{katsuyōgo}, meaning \emph{katsuyō}-words or \emph{inflectional words}, and \emph{hi-katsuyō-go} (also \emph{mu-katsuyō-go}), meaning non\emph{-katsuyō}-words or \emph{non-inflectional words}. \emph{Katsuyōgo} are verbs (\emph{dōshi} {\japanesefont 動詞}), adjectives
(\emph{keiyōshi} {\japanesefont 形容詞}) and auxiliary verbs (traditionally called verbal suffixes, \emph{jodōshi} {\japanesefont 助動詞}). On the other hand, non-inflectional words are nouns (\emph{meishi} {\japanesefont 名詞}) and adverbials (\emph{fukushi} {\japanesefont 副詞}), as well as most functional word classes, most importantly particles (\emph{joshi} {\japanesefont 助詞}).

Importantly, the term \emph{katsuyō} differentiates verbal from nominal categories. And while \emph{katusyō} is the key characteristics of verbal categories, nouns are seen as categorically uninflecting. This suggests that the term \emph{katsuyō} in fact denotes \emph{conjugation}, rather than \emph{inflection}, and we can ask ourselves whether this might not be a more fitting translation.\footnote{This is in fact the translation adopted in \citet{martin1975} for \emph{katsuyō}. On the other hand, most generative linguists use the term \emph{inflection} which is arguably the more established term. It is also found in most grammars of Japanese (see the discussion in \citealt[58--59]{uehara1998}). This discussion is reflected in \citet{uehara1998}, who uses the term `\emph{katsuyō}-inflection' as a compromise.}

A second important aspect concerns the exact dimension of \emph{inflection/katsuyō}. From an Indo-European perspective, we expect inflection to denote different morphological forms dependent on number and person as well as various TAM (temporal-aspectual-modal) categories. For example, in \ili{German}, a verb such as \emph{sprechen} `to speak' shows different forms for person and number (\emph{ich spreche} `I speak.\textsc{1sg'}, \emph{du sprichst} `you speak.\textsc{2sg'}, \emph{wir sprechen} `we speak.\textsc{1pl'}), tense (\emph{du sprachst} `you spoke.\textsc{2sg'}) and mood (\emph{du sprächest} `you would speak.\textsc{2sg.subj'}). However, Japanese lacks the categories number and person and by extension different inflectional forms for those.

Arguably, tense, aspect and modality are not expressed on the verb itself, in the first place, but via auxiliary verbs (same as, for example, \ili{German} expresses deontic modality with the auxiliary verb \emph{müssen} `must’).\footnote{With respect to tense, this was definitely the case in Classical Japanese\il{Japanese!Classical}, whereas in Modern Japanese the expression of past tense on verbs could be seen as almost fusional as it involves specific phonological processes.}

% \largerpage[2]
Rather than morphological alternation, \emph{katsuyō} denotes the conjugation of verbal categories in different inflectional forms, or \emph{katsuyōkei} ({\japanesefont 活用形}). For Modern Japanese, different authors have put forth different inventories of inflectional forms, but a standard paradigm, that was established by linguists in the Edo-period on the basis of Classical Japanese\il{Japanese!Classical}, and is nowadays still used with regard to Classical\il{Japanese!Classical}, but also Modern Japanese, contains the following six inflectional forms (see for example \citealt[222]{shibatani1990}).\footnote{When speaking of \emph{Classical Japanese}\il{Japanese!Classical} (or \emph{bungo}) I refer to the premodern written stage of the language based on the literary style of the Heian-period (794--1185) and in use till the language reform carried out during the Meiji-period (1868--1912).}

\begin{exe}
	\ex \label{12ex6}
% 	\needspace{6\baselineskip}
		I \hspace*{0.6cm} {\japanesefont 未然形} \hspace*{0.6cm} \emph{mizenkei} \hfill (negative form/irrealis) \\
		II \hspace*{0.5cm} {\japanesefont 連用形} \hspace*{0.6cm} \emph{renyōkei} \hfill (nominalization/adverbial form) \\
		III \hspace*{0.4cm} {\japanesefont 終止形} \hspace*{0.6cm} \emph{shūshikei} \hfill (sentence-final form, also conclusive) \\
		IV \hspace*{0.4cm} {\japanesefont 連体形} \hspace*{0.58cm} \emph{rentaikei} \hfill (attributive form) \\
		V \hspace*{0.54cm} {\japanesefont 已然形} \hspace*{0.56cm} \emph{izenkei} \hfill (concessive/conditional form) \\
		VI \hspace*{0.45cm} {\japanesefont 命令形} \hspace*{0.53cm} \emph{meireikei} \hfill (imperative form)\is{inflectional class!in Japanese}
\end{exe}

Some of these inflectional forms express specific functions independently, but more importantly, specific functional and lexical elements are suffixed to specific inflectional forms. For example, the negative form (\emph{mizenkei}) does not express negation by itself.\footnote{This is the case neither in Modern, nor in Classical Japanese\il{Japanese!Classical}.} Rather, it is the inflectional form an inflecting word class needs to inflect to in order to host those suffixes which negate the proposition expressed by the relevant member of this word class.

Verbs in Modern Japanese can be classified into three inflectional classes (compared to nine in Classical Japanese\il{Japanese!Classical}). Besides a group containing only two irregular verbs, which will therefore not be addressed here, inflection is especially visible for verbs belonging to the \emph{consonant inflectional paradigm}, also called \emph{consonant verbs} or \emph{consonant-stem verbs}\is{verb!consonant-stem (in Japanese)}. Those verbs end in a consonant which is followed by a vowel. It is this vowel where inflection becomes visible, as it differs across inflectional forms. There are various analyses on how to treat this vowel (see \citealt{kuno1973}; \citealt[221--235]{shibatani1990} for an overview). \citet[224]{shibatani1990} treats it as an inflectional ending, which together with the consonant root forms the stem. \citet[ch. 5]{iwasaki2013} entertains a more modern analysis of this vowel as a \emph{stem forming suffix}.

\largerpage[-1.5]

In any case, it is this vowel which shows inflection as exemplified below for the verb \emph{nom} `to drink' where typical suffixes for each inflection form are added and the alternation of the inflectional/stem forming vowel is highlighted in boldface and italics\footnote{I only gloss the vocalic elements as separate morphological elements if they allow the verb to fulfill an independent function.}.

\begin{exe}
	\ex \label{12ex7}
	\gll Watashi-ga ocha-wo nom\textbf{\emph{a}}-na-i. \\
	I-\textsc{nom} tea-\textsc{acc} drink-\textsc{neg-prs} \\
	\trans `I do not drink tea.' \hfill (negative form)
	
	\ex \label{12ex8}
	\gll Watashi-ga ocha-wo nom\textbf{\emph{i}}-ta-i. \\
	I-\textsc{nom} tea-\textsc{acc} drink-\textsc{vol-prs} \\
	\trans `I want to drink tea.' \hfill (nominalized form)
	
	\ex \label{12ex9}
	\gll Watashi-ga ocha-wo nom-\textbf{\emph{u}} koto-ga suki-da. \\
	I-\textsc{nom} tea-\textsc{acc} drink-\textsc{prs} \textsc{comp-nom} like-\textsc{cop} \\
	\trans `I like drinking tea.' \hfill (sentence-final form)
	
	\ex \label{12ex10}
	\gll Watashi-ga ocha-wo nom\textbf{\emph{e}}-ba i-i. \\
	I-\textsc{nom} tea-\textsc{acc} drink-if good-\textsc{prs} \\
	\trans `If I drank tea, it would be good.' \hfill (concessive)
	
	\ex \label{12ex11}
	\gll Ocha-wo nom-\textbf{\emph{e}}! \\
	tea-\textsc{acc} drink-\textsc{imp} \\
	\trans `Drink (the) tea!' \hfill (imperative)
\end{exe}

Exemplifying with (\ref{12ex7}), \emph{nom} `to drink' is inflected for the negative form (\emph{mizenkei}), which is expressed via the vowel \emph{-a}, yielding \emph{nom-a}, to which then the negative auxiliary verb \emph{-na} is suffixed. This auxiliary verb is itself inflecting and shares its paradigm with \emph{i-}adjectives (see \ref{12sec3.2.2}), therefore it ends in \emph{-i} in non-past. The result is \emph{nom\textbf{a}-na-i} drink-\textsc{neg-prs} `do not drink'.

In contrast, the sentence-final and the attributive forms, which have become unified in Modern Japanese for all word classes but \emph{na}-adjectives -- and thus for verbs can no longer be differentiated \citep{kuno1973, frellesvig2010}, which is why they are not given separately -- fulfill independent functions as does the imperative form.

This classification in six inflection forms is not shared by all authors (see \citealt[229]{shibatani1990}). For example, \citet[27--28]{kuno1973} has proposed different inflectional paradigms on the basis of Modern Japanese, illustrated below for \emph{hanas} `to speak'.

\begin{table}[H]
	\centering
	\begin{tabularx}{0.8\textwidth}{lll}
		\lsptoprule
			Inflection Form & Verb & Translation \\
			\midrule
			Present/Non-Past & hanas-u & `speak’ \\
			Past & hanash-ita & `spoke’ \\
			Imperative & hanas-e & `speak!’ \\
			Cohortative & hanas-ō & `let’s speak’ \\
			Continuative & hanash-i & `speaking’ \\
			Gerundive & hanash-ite & `speak-and’ \\
			Conditional & hanas-eba & `if…speak’ \\
			Perfect Conditional & hanash-itara & `if…spoken’ \\
			Perfect Suppositional & hanash-itarō & `I suppose…spoke’ \is{inflectional class!in Japanese}\is{verb!consonant-stem (in Japanese)}\\
		\lspbottomrule
	\end{tabularx}
	\caption{Inflectional forms of Modern Japanese according to \citet{kuno1973}}
	\label{12tab1}
\end{table}

On the other hand, inflection is harder to see for so-called \emph{vowel-stem verbs}\is{verb!vowel-stem (in Japanese)}, or \emph{vowel verbs}. This is because their stem already ends in a vowel, i.e., according to \citet[80]{iwasaki2013}, their stem forming suffix is Ø. Nevertheless, they are typically still counted as an inflecting category (\citealt{shibatani1990}; \citealt{uehara1998}; \citealt{iwasaki2013} among others). Their paradigm is shown below for the vowel-stem verb \emph{tabe} `to eat'.

\begin{exe}
	\ex \label{12ex12}
	\gll Watashi-ga gohan-wo tabe-na-i. \\
	I-\textsc{nom} rice-\textsc{acc} eat-\textsc{neg-prs} \\
	\trans `I do not eat rice.' \hfill (negative form)
	
	\ex \label{12ex13}
	\gll Watashi-ga gohan-wo tabe-ta-i. \\
	I-\textsc{nom} rice-\textsc{acc} eat-\textsc{vol-prs} \\
	\trans `I want to eat rice.' \hfill (nominalized form)
	
	\ex \label{12ex14}
	\gll Watashi-ga gohan-wo tabe-ru koto-ga suki-da. \\
	I-\textsc{nom} rice-\textsc{acc} eat-\textsc{prs} \textsc{comp-nom} like-\textsc{cop} \\
	\trans `I like eating rice.' \hfill (sentence-final form)
	
	\ex \label{12ex15}
	\gll Watashi-ga gohan-wo tabe-reba i-i. \\
	I-\textsc{nom} rice-\textsc{acc} eat-if good-\textsc{prs} \\
	\trans `If I ate rice, it would be good.' \hfill (concessive)
	
	\ex \label{12ex16}
	\gll Gohan-wo tabe-ro! \\
	rice-\textsc{acc} eat-\textsc{imp} \\
	\trans `Eat (the) rice!' \hfill (imperative)
\end{exe}

\subsubsection{Inflection of adjectives\is{adjective!Japanese|(}} \label{12sec3.2.2}

Japanese has two canonical adjective groups, one referred to here as \emph{i}-adjectives, the other as \emph{na}-adjectives. Whether \emph{na}-adjectives actually inflect is subject to debate. In fact, this debate is tied to the very understanding of \emph{katsuyō}\is{katsuyō|)} (cf. \citealt[59--64]{uehara1998}; \citealt{kishimotouehara2016}). The reason why \emph{i}-adjectives are usually perceived as an inflecting category is illustrated in (\ref{12ex17}) through (\ref{12ex20}) below in which different inflection forms and, if applicable, typical suffixes are given for the \emph{i}-adjective \emph{aka-i} `red'.

\begin{exe}
	\ex \label{12ex17}
	\gll Kono hana-wa aka-\textbf{\emph{i}}. \\
	this flower-\textsc{top} red-\textsc{prs} \\
	\trans ‘This flower is red.’ \hfill (sentence-final form)
	
	\ex \label{12ex18}
	\gll Kono hana-wa aka\textbf{\emph{ku}}-na-i. \\
	this flower-\textsc{top} red-\textsc{neg-prs} \\
	\trans ‘This flower is not red.’ \hfill (negative form)
	
	\ex \label{12ex19}
	\gll hana-wa aka\textbf{\emph{ke}}-reba \\
	flower-\textsc{top} red-if \\
	\trans ‘if the flower was red’ \hfill (concessive)
	
	\ex \label{12ex20}
	\gll Kono hana-wa aka\textbf{\emph{k}}-atta. \\
	this flower-\textsc{top} red-\textsc{pst} \\
	\trans ‘This flower was red.’ \hfill (past)
\end{exe}

At first glance, \emph{na}-adjectives have to co-occur with a copula which is itself inflecting. The respective patterns are given below for the \emph{na}-adjective \emph{nigiyaka-na} `lively'.\footnote{For reasons of space, the question of which adjective groups in Japanese inflect will not be discussed here. Readers are referred to \citet{uehara1998}, \citet{nishiyama1999}, \citet{yamakido2005} among others.}

\begin{exe}
	\ex \label{12ex21}
	\gll Kono machi-wa nigiyaka-da. \\
	this city-\textsc{top} lively-\textsc{cop.prs} \\
	\trans ‘This city is lively.’ \hfill (sentence-final form)
	
	\ex \label{12ex22}
	\gll Kono machi-wa nigiyaka-de-wa-na-i. \\
	this city-\textsc{top} lively-\textsc{cop-top-neg-prs} \\
	\trans ‘The city is not lively.’ \hfill (negative form)
	
	\ex \label{12ex23}
	\gll machi-wa nigiyaka-de-areba \\
	city-\textsc{top} lively-\textsc{cop}-if \\
	\trans ‘if the city was lively’ \hfill (concessive)
	
	\ex \label{12ex24}
	\gll Kono machi-wa nigiyaka-datta. \\
	this city-\textsc{top} lively-\textsc{cop.pst} \\
	\trans ‘This city was lively.’ \hfill (past)\is{adjective!Japanese|)}
\end{exe}

\subsubsection{Non-inflection of nouns\is{uninflecting!word class|(}\is{noun!Japanese|(}} \label{12sec3.2.3}

As mentioned above, the other major lexical word class, the nominal class, is uninflecting. Japanese lacks number also in the nominal domain, but it does have case. However, case is not expressed via inflection but via case particles which are usually treated as (themselves noninflecting) clitics \citep{uehara1998, iwasaki2013}. This is exemplified below.

\begin{exe}
	\ex \label{12ex25}
	\gll Sake-\textbf{\emph{ga}} suki-da. \\
	alcohol-\textsc{nom} fond-\textsc{cop.prs} \\
	\trans ‘I like alcohol.’ \hfill (nominative)
	
	\ex \label{12ex26}
	\gll Watashi-ga sake-\textbf{\emph{wo}} nom-u koto-ga suki-da. \\
	I-\textsc{nom} alcohol-\textsc{acc} drink-\textsc{prs} \textsc{comp-nom} fond-\textsc{cop.prs} \\
	\trans ‘I like drinking alcohol.’ \hfill (accusative)
	
	\ex \label{12ex27}
	\gll sake-\textbf{\emph{ni}} mizu-wo ire-ru \\
	alcohol-\textsc{dat} water-\textsc{acc} put.in-\textsc{prs} \\
	\trans ‘to put water in alcohol’ \hfill (dative/locative)
\end{exe}

This discussion emphasizes that the existence of \emph{uninflectable}\is{uninflectable!element} lexemes is not expected in Japanese in the system presented here. There should be no lexeme that belongs to an inflecting word class, but fails to show the relevant morphology associated with this word class\is{noun!Japanese|)}.\footnote{Potentially, this might be related to the fact that uninflectable lexemes might be more often found in the realm of nominals than verbals in general \citep[143]{spencer2020}. Nevertheless, we do find uninflectable verbs in some languages as demonstrated by Spencer (this volume).}

\subsection{\emph{Rentaishi}} \label{12sec3.3}

\subsubsection{Their development} \label{12sec3.3.1}

The question is then how \emph{rentaishi} fit into the picture. This word class literally translates to `adnominal-noun-word' and is in \ili{English} usually referred to as \emph{adnoun} or \emph{prenoun}, \ili{German} equivalents are \emph{Adnomen}, sometimes \emph{Adnominalie} \citep{rickmeyer1995}. An alternative name is \emph{fukutaishi}, which means something like `added word'. \emph{Rentaishi} can be described as a word class with very few lexemes that have all been derived from other word classes and have become fossilized. Syntactically, they are restricted to their role as attributive modifiers. Morphologically, they can be defined as being rather heterogeneous which also sets them apart from the word classes of verbs and adjectives in Japanese which, as seen above, are normally characterized via their distinct morphology (cf. \citealt{nittamurakishibataniyazawa2000}).

The origin of \emph{rentaishi} can be traced back to the Meiji-period (1869--1912), thus they entered the Japanese word class system later than the other major word classes \citep{frellesvig2010}. The first mention of this word class can be attributed to MITSUYA Shigematsu in his lectures given in 1915/16, but published in 1932. Mitsuya describes \emph{rentaishi} as ``attributives that are entirely uninflecting'', that are ``adjective-like without being adjectives'' \citep[238]{mitsuya1932}. The important linguists of the time, especially MATSUSHITA Daizaburō (1878--1935), realized that there are several lexemes in Japanese that are restricted to attributive use, do not inflect and whose sole function is to modify nouns. In his \emph{Revised Standard Grammar} (\emph{Kaisen hyōjun Nihon bunpō}), \citeauthor{matsushita1976} (\citeyear{matsushita1976} [1928]), who calls them \emph{fukutaishi}, writes.

\begin{exe}
	\ex \label{12ex28}
	Fukutaishi (Adjective) are lexemes that express an attributive concept which is in turn dependent on another concept, that modify the meaning of other lexemes, and that do not have a predicative character.
	
	\hfill (\citeauthor{matsushita1976} \citeyear{matsushita1976} [1928]: 204, [bracket in original])
\end{exe}

This highlights the following four characteristics of \emph{rentaishi}.

\begin{exe}
	\ex \label{12ex29} Characteristics of \emph{rentaishi} according to \citeauthor{matsushita1976} (\citeyear{matsushita1976} [1928]):
	\begin{enumerate}
		\item lexical status
		\item dependency
		\item modifiers
		\item no predicative use
	\end{enumerate}
\end{exe}

Due to the influence of Matsushita and their incorporation into \citeauthor{hashimoto1948}'s Standard Grammar (\citeyear{hashimoto1948} [1934]), in which Hashimoto highlights their lexemic status as well and puts special emphasis on their role as attributive modifiers of nominals and their inability to appear as subjects of a sentence (\citealt{hashimoto1948} [1934]: 63--65), \emph{rentaishi} had become a recognized member of the Japanese word class system.

Although neither author emphasizes the characteristic \emph{uninflectability}, we will see that the perception of \emph{rentaishi} as an independent category is very much tied to this notion.

\subsubsection{Syntactic characteristics and morphological subgroups} \label{12sec.3.3.2}

Potentially, the most obvious cases of \emph{rentaishi} are lexemes that are derived from verbs and have become fossilized in attributive form. As briefly mentioned above, the sentence-final form (\emph{shūshikei}) and the attributive form (\emph{rentaikei}), cannot be differentiated for verbs in Modern Japanese, but used to be different in Classical Japanese\il{Japanese!Classical}.\footnote{This unification potentially happened in Late Middle Japanese (12--16 AD; \citealt{lewin2003, frellesvig2010}).} While the sentence-final form of verbs belonging to the \emph{ri}-inflectional-paradigm ended in \emph{-ri}, the attributive form ended in \emph{-ru}.\footnote{In \citeauthor{iwasaki2013}'s terms (\citeyear{iwasaki2013}), this would mean that the stem forming suffix of the \emph{shūshikei} was \emph{-i}, that of the \emph{rentaikei} \emph{-u}.} Verbal derivates belonging to the group of \emph{rentaishi} are solely derived from verbs of the \emph{ri}-paradigm, thus end in \emph{-ru}. Examples include the following.\footnote{Note that \emph{aru} `(a) certain` is homophonous to the existential verb/copula \emph{aru} `to exist` which is why some	authors treat them as the same element (\citealt{lehmannnishina2015} a.o.).  However, the existential verb \emph{aru} is restricted to inanimate entities and with these entities can express predication	relations, which is why these two elements are kept apart here since the \emph{rentaishi aru} never appears in predicative position.}

\begin{exe}
	\ex \label{12ex30} 
	\emph{aru} `(a) certain' \\
	\emph{saru} `previous' \\
	\emph{akuru} `following' \\
	\emph{kitaru} `ahead'
\end{exe}

All those lexemes are impossible in predicative position, illustrated with \emph{akuru}.

\begin{exe}
	\ex \label{12ex31}
	\begin{xlist}
		\ex 
		\gll akuru toshi \\
		following year \\
		\trans ‘the following year’
		
		\ex[*] {
			\gll Toshi-wa akuru. \\
			year-\textsc{top} following \\
			\trans Intend.: ‘The year is the following.’}
	\end{xlist}
\end{exe}

There are other verbal derivates in which a lexemic verb is first combined with the suffix \emph{-yu} (or \emph{-iu)}, which was used to express passive and spontaneity in Classical Japanese\il{Japanese!Classical}. This suffix is then itself inflected for attributive form by means of the suffix \emph{-ru}. Representative examples are \emph{iwayuru} `so-called' and \emph{arayuru} `all kinds of'. Exemplifying with \emph{iwayuru}, this word can be divided into the verb \emph{iu} `to say', which is inflected for \emph{mizenkei}, (negative form) and the marker for spontaneity/passive \emph{-yu}. This marker is in turn inflected for the attributive form. Again, the \emph{rentaishi} \emph{iwayuru} is not used in Modern Japanese, except for attributive contexts.

\begin{exe}
	\ex \label{12ex32}
	\begin{xlist}
		\ex 
		\gll iwayuru gakusei \\
		so-called student \\
		\trans ‘a so-called student’
		
		\ex[*] {
			\gll Kono gakusei-wa iwayuru. \\
			this student-\textsc{top} so-called \\
			\trans Intend.: ‘This student is so-called.’}
	\end{xlist}
\end{exe}

There are also members of the word class \emph{rentaishi} which are are derived from adjectives. Besides those ending in the element \emph{-na} to be discussed in Section \ref{12sec5}, there are those which co-occur with the attributive form of the premodern copula \emph{nari} or \emph{tari}.\footnote{\emph{tari} can be seen as an allomorph which was restricted to Sino-Japanese lexemes in Classical Japanese\il{Japanese!Classical}.} Again, the relevant lexemes in Modern Japanese are restricted to mono-attributive use.

\begin{exe}
	\ex \label{12ex33}
	\begin{xlist}
		\ex 
		\gll tannaru machigai \\
		mere misunderstanding \\
		\trans ‘a mere misunderstanding’
		
		\ex[*] {
			\gll Machigai-wa tannaru. \\
			misunderstanding-\textsc{top} mere \\
			\trans Intend.: ‘The misunderstanding is mere.’}
	\end{xlist}
\end{exe}

Finally, there are \emph{rentaishi} derived from nouns. An example is \emph{waga} `my, our' which is derived from \emph{wa}, the noun used in Classical Japanese\il{Japanese!Classical} to refer to the 1\textsuperscript{st} person singular, and \emph{-ga}, which is the nominative particle, but was also used as the genitive particle in Old, and partly in Classical Japanese\il{Japanese!Classical} \citep{frellesvig2010}. \emph{Waga} is basically used only in some fixed expression in Modern Japanese, such as in (\ref{12ex34}) given below. The same restrictions apply.

\begin{exe}
	\ex \label{12ex34}
	\begin{xlist}
		\ex 
		\gll waga kuni \\
		our country \\
		\trans ‘our country’
		
		\ex[*] {
			\gll Kono kuni-wa waga. \\
			this country-\textsc{top} our \\
			\trans Intend.: ‘This country is ours.’}
	\end{xlist}
\end{exe}

There are other lexemes which can be classified as nominal \emph{rentaishi}. Those end in \emph{no}, the particle used to mark genitive and a variation of other adnominal functions in Modern Japanese (for example \citealt{saitolinmurasugi2008}). In Classical Japanese\il{Japanese!Classical}, this particle was used also to mark the subject, although some uses are attested where it attaches to a nominal modifier within the DP. Examples include \emph{reino} (or potentially \emph{rei-no}) `ordinary, aforementioned' and \emph{honno} (\emph{hon-no}) `mere'. Neither do these lexemes themselves nor the part preceding \emph{no} occur predicatively. We will get back to this group in Section \ref{12sec5}.

\begin{exe}
	\ex \label{12ex35}
	\begin{xlist}
		\ex 
		\gll honno kodomo \\
		mere kid \\
		\trans ‘mere kid’ \hfill \citep[122]{nagano2016}
		
		\ex[*] {
			\gll Kono kodomo-wa honno (da). \\
			this kid-\textsc{top} mere (\textsc{cop}) \\
			\trans Intend.: ‘This kid is mere.’}
	\end{xlist}
\end{exe}

Table \ref{12tab2} illustrates all relevant subgroups of \emph{rentaishi} (cf. \citealt[\mbox{181--182}]{kohlich2020})\is{uninflecting!word class|)}.

\begin{table}[ht]
	\centering
	\begin{tabularx}{\textwidth}{l>{\raggedright}p{4.7cm}Q}%lL{4.7cm}L{4cm}
		\lsptoprule
			Group & Subgroups & Examples \\
			\midrule
			verbal \emph{rentaishi} & attributive form of verbs in non-past & \emph{aru} ‘(a) certain’; \mbox{\emph{saru} ‘previous’}; \mbox{\emph{akuru} ‘following’}; \mbox{\emph{kitaru} ‘ahead’} \\[0.6cm]
			& attributive form of verbs with morpheme \emph{-iu} & \emph{iwayuru} ‘so-called’; \emph{arayuru} ‘all kinds of ’ \\[0.6cm]
			& attributive form of verbs in past & \emph{taishita} ‘considerable’; \emph{tonda} ‘remarkable’ \\[0.6cm]
			adjectival \emph{rentaishi} & lexemes ending in attributive forms of premodern copula \emph{-naru/taru} & \emph{tannaru} ‘mere’; \mbox{\emph{shutaru} ‘main’} \\[0.6cm]
			& three lexemes in combination with element \emph{na} & \emph{ōkina} ‘big’; \mbox{\emph{chiisana} ‘small’}; \mbox{\emph{okashina} ‘strange’} \\[0.6cm]
			nominal \emph{rentaishi} & lexemes ending in \emph{-ga}, often in particular collocations & \emph{waga} ‘my/our’; \mbox{\emph{onoga} ‘my’} \\[0.6cm]
			& lexemes ending in \emph{-no} & \emph{reino} ‘ordinary, aforementioned’; \mbox{\emph{honno} ‘only, mere’} \\
		\lspbottomrule
	\end{tabularx}
	\caption{Different morphological subgroups of \emph{rentaishi}}
	\label{12tab2}
\end{table}

\section{The historical development of \emph{rentaishi}} \label{12sec4}

As emphasized by \citeauthor{corbett2025} (this volume, p. 43) ``loss of inflection can reflect a change of part of speech''. Having outlined the syntactic restrictions and morphological heterogeneity of \emph{rentaishi}, we are in a position to analyze the relevance of this word class for the debate on \isi{uninflectedness}. This will be done by taking the historic circumstances of its development into consideration, specifically in connection to the Meiji-period, the field of \emph{Dutch Studies}\il{Dutch}\is{Dutch Studies} and their connection to adjectives.

\subsection{The Meiji-period} \label{12sec4.1}

The word class \emph{rentaishi} was established and subsequently became the object of interest towards the end of the Meiji-period, thus around the beginning of the 19\textsuperscript{th} century. Historically, this period is characterized by a stronger interest of Japan in the West.
After decades of very limited contact, the intellectuals of the Meiji-period sought out the Western cultural and political achievements and started comparing them to those of their own country. This trend also pertained to the field of linguistics (cf. \citealt{heinrich2002} among others). Several linguists started turning their eyes towards Western linguistics, comparing the languages of the West with the Japanese language as well as linguistic traditions in all areas, including grammar, but also word classes (cf. \citealt[47--54]{sugimoto1991}).\footnote{Important linguists were TANAKA Yoshikado (1841--1879), NAKANE Kiyoshi (NAKANE Kōtei, 1839--1913) and NAKAGANE Masahira (1824--1884).} In turn, this also led to a better understanding of the Japanese language.

Another important trend of the Meiji-Period was the beginning unification of the Japanese language, known as \emph{\isi{genbun'itchi}}. Residues of diglossia between the written language, which mostly followed the norms of Classical Japanese\il{Japanese!Classical}, and the spoken language had hindered a unification of the Japanese language for decades, but from around 1885 onward, as a result of the \emph{\isi{genbun'itchi}} movement, a standardized unified language was gradually achieved. Notably, this trend was tied to the fact that Japanese linguists started to realize that this lack of unification was not observable in other countries \citep{gottlieb1991}.

\subsection{Adjectives in \emph{Dutch Studies}\il{Dutch|(}\is{Dutch Studies|(}} \label{12sec4.2}

The two languages of interest during the Meiji-period were \ili{English}, which was identified as the most important language of the West, and Dutch. Linguists were interested in Dutch since it had been the most important language of trade during the Edo-period when trade with outsiders was formally forbidden and was solely carried out with the Netherlands in the artificial island of Dejima, close to Nagasaki \citep{sugimoto1967, sugimoto1991, nespoli2023}. The study of Dutch had in fact already started during the Edo-period and resulted in the education of several translators and interpreters between the two languages as well as several translations of important Dutch works in various areas, including linguistics. This study of the Dutch language in Japan is referred to as \emph{Dutch Studies} or \emph{rangaku} {\japanesefont 蘭学} (ibid.).\footnote{As argued in \citet[48]{nespoli2023}, it might be more precise to refer to the dedicated studies of the Dutch language as \emph{rangogaku}. However, \emph{rangaku} was the preferred term used by scholars of the time as also evidenced by the titles of books and studies concerning the Dutch language.}

The emergence of \emph{rentaishi} as a word class, perceived independently by Japanese linguists, is tied closely to \emph{rangaku}. This is because the linguists of the Meiji-period came to realize that the Japanese adjectives, or \emph{keiyōshi}, are an inflecting category, and more importantly a category that inflects \emph{independently}. In their view, this contrasted them sharply from the adjectival category in Dutch, and by extension all European languages as they thought, which only inflects in accordance with the head noun and never appears in predicative position without the aid of the copula \citep{kim2006}.

In this respect, note that also in the European tradition, the adjective had for a long time been regarded as \emph{\isi{nominem adjectivum}} `added noun', following the first terminology of Thomas of Erfurt\ia{Thomas of Erfurt} (12\textsuperscript{th} century), whose purpose it is to follow the \emph{\isi{nominem subjectivum}} and to acquire its features \citep{hajek2004}. Dutch linguists of the 18\textsuperscript{th} century, noteworthy all those relevant to the present article as they were perceived by contemporary Japanese linguists,\footnote{This concerns essentially works of Willem Séwel\ia{Séwel, Willem} (1653 -- 1720), François Halma\ia{Halma, François} (1653 -- 1722) and Pierre Marin\ia{Marin, Pierre} (1667/8 -- 1718) (cf. \citealt{nespoli2023}).} referred to the Dutch adjective\is{adjective!Dutch} as \emph{bijvoeglijk naamwoord} or \emph{toevoeglijk naamword}, both approximately translatable as `added noun' (\citealt[624]{sugimoto1991}; \citealt[128]{kim2006}; \citealt{nespoli2023}).\footnote{Compare also the \ili{German} terms \emph{beystendiges Nennwort} and \emph{Beynahmen}, close to `accompanying noun`, used by \ili{German} linguists in the 18\textsuperscript{th} century \citep[1--5]{spitzl-dupic2011}.} What is reflected in these terms is that adjectives do not occur independently, but only together with an \emph{independent noun}, called among others \emph{zelfstandig naamwoord} (ibid.). Crucially, adjectives are dependent because they need to inflect according to the \emph{phi}-features exhibited by the head noun they modify.\footnote{This only concerns their use in attributive position, whereas they do need a copula in order to appear in predicative position and show no inflection. See next section.}

The Japanese linguists working on Dutch and writing their own Dutch grammars created various Japanese translations, or rather equivalents, for \emph{bijvoeglijk} (\emph{toevoegelijk}) \emph{naamwoord}. Some of these terms were neologisms \citep{kim2006, sugimoto1991}, but others were calques taken either from the tradition of \emph{kanbunten} (`Grammar of Chinese', cf. \citealt[44--45]{sugimoto1991}) or earlier works in the tradition of \emph{kokugogaku} (for example from OGYŪ Sorai (1666--1728)\ia{Ogyū, Sorai}, \citealt[395--399]{nespoli2023}). A (inconclusive) list of relevant terms and in which works they are given is presented below (cf. \citealt[188]{kohlich2020}; especially the appendix in \citealt{sugimoto1991}).

\begin{table}[ht]
	\centering
	\begin{tabularx}{\textwidth}{QQQ}
		\lsptoprule
			Author & Book & Adjective \\
			\midrule
			Nakano Ryūho (Shizuki Tadao, Wilgen Akker, 1760–1806)\ia{Nakano, Ryūho} & \emph{Ryūho Nakano sensei bunpō} (Grammar of Master Nakano Ryūho, around 1798) & \emph{bōki-myōgo} {\japanesefont 傍寄名語} (lit. ‘noun depending on what is nearby’) \\
			Baba Sadayoshi (Baba Sajūrō, 1787–1822)\ia{Baba, Sadayoshi} & \emph{Teisei rango kyūhin-shū} (Collection of the nine parts of speech of Dutch, revised, 1814)\footnotemark & \emph{kyoseishi} {\japanesefont 虚静詞} (\emph{kyo-} = empty, \emph{-sei-} = stative, \emph{-shi} = word) \\
			Yoshio Shunzō\ia{Yoshio, Shunzō} (1787–1843) & \emph{Rokkaku zenpen} (On the six cases, first edition, 1814) & \emph{baimeishi} {\japanesefont 陪名詞} (lit. `added/dependent noun') \\
			Fujibayashi Taisuke (Fujibayashi Fuzan, 1781–1836)\ia{Fujibayashi, Taisuke} & \emph{Oranda gohō kai} (Analysis and grammar of Dutch, 1815) & \emph{fuzokumeigen} {\japanesefont 附屬名言} (`added noun') \\
			Ōtsuki Genkan (Ōtsuki Banri, 1785–1837)\ia{Ōtsuki, Genkan} & \emph{Rangaku-han} (Dutch colloquial speech, 1816, 1824) & \emph{setsumeishi} {\japanesefont 接名詞} (`connected noun') \\
			Ōba Sessai (1805–1873)\ia{Ōba, Sessai} & \emph{Yaku Oranda bungo} (Dutch written language, 1855) & \emph{baiji}\is{baiji} {\japanesefont 陪辭} (`added/dependent word') \\
		\lspbottomrule
	\end{tabularx}
	\caption{Different terms for adjectives in \emph{Dutch} grammars in Japan}
	\label{12tab3}
\end{table}

\footnotetext{Presumably, this is a revised version of \emph{Oranda-go shihinkō} {\japanesefont 和蘭語詞品考} (Thoughts on the parts of speech of Dutch) by Nakano Ryūho\ia{Nakano, Ryūho}. The manuscript of this title cannot be identified, although some authors believed that this work is the same as \emph{Ryūho Nakano sensei bunpō}\ia{Nakano, Ryūho}. The fact that the terms for the Dutch adjective\is{adjective!Dutch} are different in the two works, for example, does suggest otherwise. Cf. \citet{sugimoto1991} and \citet[394]{nespoli2023} for discussion.}

All authors characterize the adjective as a lexeme, via the suffixes \emph{-shi} or \emph{-gen}. Some authors are more explicit and refer to them as \emph{nouns} (\emph{meigen} and \emph{meishi}), while others do not. What all these terms reflect is the \emph{added} and \emph{dependent} character of the adjective. This is evident in the prefixes \emph{fuzoku}- (`affiliation'\emph{,} `membership'), \emph{setsu}- (`connection') and \emph{bai}- (`addition'). The morpheme \emph{sei} `stative' differentiates adjectives from \emph{dynamic} (\emph{dō}) verbs\il{Dutch|)}\is{Dutch Studies|)}.

\subsection{Adjectives and	\emph{rentaishi}\is{uninflecting!word class|(}} \label{12sec4.3}

At the same time, the linguists of the Meiji-period realized that the Japanese adjectives, or \emph{keiyōshi} `quality words', can inflect independently for tense, in attributive as well as in predicative position, and do not show agreement with their head noun. That is, they do not simply just mirror the (phi-)features of nouns and therefore are not \emph{added} or \emph{accompanying}. In accordance with the recent popularization of linguistic comparisons of Japanese and Western languages, adjectives were stylized as \emph{independent} against the Western adjectives, \emph{dependent} and \emph{added}. This view, and the historical development depicted on the preceding pages, is reflected in the works of ŌTSUKI Fumihiko (1847--1928), the grandson of ŌTSUKI Gentaku (1757--1827)\ia{Ōtsuki, Gentaku}, one of the most important experts in \emph{Dutch Studies}\il{Dutch}\is{Dutch Studies}, and nephew of ŌTSUKI Genkan\ia{Ōtsuki, Genkan}. Ōtsuki Fumihiko wrote two of the first relevant grammars of Japanese as it approached its modern stage: \emph{Genkai} (Sea of Words, \citeyear{otsuki1891}) and \emph{Kō-Nihon bunten} (Detailed Japanese Grammar, \citeyear{otsuki1897}). In \emph{Genkai}, Ōtsuki writes the following.

\largerpage[-0.5]

\begin{exe}
	\ex \label{12ex36}
	The \ili{English} \emph{Adjective} [sic!] merely precedes the noun and bestows it with qualitative features. The Japanese adjective [keiyōshi, V.K.] also adds qualitative features to the noun, however in terms of morphological build-up, it is entirely different. The fact that it shows inflection in word-final position and different inflection forms is a feature it shares with verbs. Moreover, placed after the noun, it marks the end of a sentence.
	
	\hfill \citep[28]{otsuki1891}
\end{exe}

In this definition, it becomes obvious that Ōtsuki sees the Japanese adjective more powerful than the \ili{English}, the Western, adjective. This is illustrated below with the examples repeated from above where special emphasis should be placed on the \ili{English} translation. Although this is not the precise state of the language during Ōtsuki's time, the examples show that the Japanese adjective occurs independently, whereas the \ili{English} adjective is always accompanied by a copula when it appears in predicative position.\footnote{Note that I give an example of an \emph{i}-adjective here, abstracting away from the question in how far \emph{na}-adjectives show inflection. In fact, all adjectives, also those inflected for past or bestowed with the negative suffix, are possible in attributive position without additional elements such as relative clause markers.}

\begin{exe}
	\ex \label{12ex37}
	\begin{xlist}
		\ex 
		\gll Hana-wa aka-i. \\
		flower-\textsc{top} red-\textsc{prs} \\
		\trans ‘The flower is red.’
		
		\ex 
		\gll Hana-wa akaku-nai. \\
		flower-\textsc{top} red-\textsc{neg} \\
		\trans ‘The flower is not red.’
		
		\ex 
		\gll Hana-wa aka-katta. \\
		flower-\textsc{top} red-\textsc{pst} \\
		\trans ‘The flower was red.’
	\end{xlist}
\end{exe}

{\sloppy YAMADA Yoshio (1873--1958), the linguistic figure following Ōtsuki, argues that in Japanese, both verbs and adjectives\is{adjective!Japanese}, have the ``ability of predication'' (\emph{``chinjutsu-no nōryoku''}), an ability which should only apply to inflecting word classes \citep[861]{yamada1908}. In contrast, he describes \ili{English} adjectives as follows.\par}

\begin{exe}
	\ex \label{12ex38}
	The \ili{English} \emph{keiyōshi}, in other words the \emph{adjective} […], should be considered a \emph{baiji}\is{baiji}, but not a \emph{keiyōshi}. 
	
	\hfill \citep[869]{yamada1908}
\end{exe}

Initially, Yamada uses the term \emph{keiyōshi} with respect to both languages, \ili{English} and Japanese. Afterwards, he decides to refer to the \ili{English} adjective as \emph{baiji}\is{baiji} or `added word', a direct relict of the \emph{rangaku} terms \emph{baiji}\is{baiji} used by ŌBA Sessai\ia{Ōba, Sessai} and \emph{baimeishi} used by YOSHIO Shunzo\ia{Yoshio, Shunzō}. Yamada's use of the term \emph{baiji}\is{baiji} can be interpreted as pejorative and possibly even politically motivated, highlighting the superiority of the Japanese \emph{keiyōshi} \citep{kim2006}.

Having established the superiority and the more independent character of the Japanese, contra the Western, adjective, Japanese linguists also became aware that there is a group of attributive modifiers in Japanese which do not inflect, thus do not appear independently at the end of a sentence and express tense. One such element is \emph{aru} `a certain'. They realized that such elements could hardly be \emph{keiyōshi} as they are not independent. Furthermore, a lexeme such as \emph{aru,} which looks verbal on the morphological level, does not share any features with verbs as it lacks all inflectional forms of the verbal paradigm. It therefore could also hardly be classified as a verb.

Building on the grammatical descriptions of \citet{otsuki1891} and \citet{yamada1908}, \citeauthor{matsushita1976} (\citeyear{matsushita1976} [1938]) introduced \emph{rentaishi} (\emph{fukutaishi}) in his grammar and saw several successors (\citealt{sekine1937, kieda1938} among others). Those linguists recognized that many lexemes in Japanese simply do not have any other function than to modify nouns, and that most of these lexemes are fossilized derivates of other word classes and frozen in their morphological form and usage domain.

Keeping apart the \ili{English} \emph{adjective} from the Japanese \emph{keiyōshi}, I argue (and see \citealt{kim2006} for an analysis similar in spirit) that the category of \emph{rentaishi} was created subsequently in order to fill the slot of dependent attributive modifiers in Japanese. All those attributive modifiers derived from other word classes and unable to inflect were sorted into this newly formed word class. Thus, Japanese linguists avoided the existence of \emph{uninflectable} lexemes\is{uninflectable!element|(}, maintaining the variable \emph{inflection}, or \emph{katsuyō}\is{katsuyō}, as the main difference between verbal and nominal categories in Japanese and predicting that word classes are either entirely inflecting or entirely uninflecting. This is why authors such as Mitsuya describe \emph{rentaishi} as ``adjective-like without being adjectives'' \citep[238]{mitsuya1932}. In a nutshell, the integration of the word class \emph{rentaishi} marked an important turning point in linguistics in Japan, highlighting the criterion \emph{inflection} and setting Japanese word classes in relation to other languages.

From the perspective of (un-)inflectedness, the ability to inflect is clearly connected to word class status. Uninflectable lexemes were not part of the Japanese word class system during the Meiji-period, and presumably until the middle of the 20\textsuperscript{th} century. Rather than accepting the existence of uninflectable lexemes, Japanese linguists established a new word class which is characterized via the very feature [+uninflecting].

Finally, consider what \citet{baermanbrowncorbett2005} call \emph{neutralization}\is{neutralisation} and define as follows:

\begin{exe}
	\ex \label{12ex39}
	\begin{xlist}
		\exi{(i)} In the presence of a particular combination of values of one or more other features (the context), there is a general loss of all values of a particular feature F found elsewhere in the language.
		\exi{(ii)} No syntactic objects distinguish any values of feature F in the given context, and feature F is therefore syntactically irrelevant in that context.
	\end{xlist}
	\hfill \citep[30]{baermanbrowncorbett2005}
\end{exe}

Neutralization\is{neutralisation} basically accounts for cases in which a feature becomes syntactically irrelevant in certain contexts. The example given by the authors is \emph{gender} in plural contexts in \ili{Russian} \citep[28]{baermanbrowncorbett2005} as there is simply no morphological distinction on nouns according to this feature in this context. This means that, different to \isi{uninflectedness}, neutralization\is{neutralisation} presupposes the existence of a specific feature that is relevant in syntax.

But that is not the case for \emph{rentaishi}, at least not as depicted here, since no inflection is observable in the first place. A case of \emph{neutralization}\is{neutralisation} can indeed be put forth for in-between stages especially of verbal \emph{rentaishi} in which they were still perceived as verbs but only used in particular contexts, that is, in attributive position. \citet[356--359]{kieda1938} still recognized an in-between stage for lexemes such as \emph{aru} `(a) certain' which is why he decides not to classify them as \emph{rentaishi} \citep[359]{kieda1938}. This means we manage to depict an evolution, at least according to the view of this author, from a neutralized\is{neutralisation} lexeme to a member of an uninflecting word class\is{uninflecting!word class|)}.

\section{Discussion: The modern dimension of \emph{rentaishi}} \label{12sec5}

Over the years, the understanding and acceptance of \emph{rentaishi} as a word class, predominantly in Japan, but also abroad, has shifted. \citet[53]{komatsu1973} documents a rise in critical studies and surveys of this category and its inclusion into the Japanese word class system, following their establishment towards the middle of the 20\textsuperscript{th} century. These critical studies include \citet{tokuda1934} and \citet{suzuki1959} and prominently \citet{watanabe1971} and were followed by a gradual marginalization of this category. Nowadays, \emph{rentaishi} are absent from most Standard Grammars (\citealt{hinds1988, tsujimura2014, hasegawa2015} among others).

While some authors have performed dedicated analyses of \emph{rentaishi},\footnote{For example, \citet{watanabe1971} treats \emph{rentaishi} as \emph{adnominal auxiliaries}	(\emph{rentai-fukushi}), \citet[742--753]{martin1975} treats them as defective nouns and \citet[27]{kato2003} as a subgroup of dependent nominals.} most authors who explicitly argue against the word class \emph{rentaishi} instead sort relevant lexemes into the groups they morphologically belong to (\citealt{mio1942, mio1958, suzuki1959, lehmannnishina2015}). That means a lexeme such as \emph{aru} `(a) certain' is classified as a verb, a lexeme such as \emph{waga} `my, our' as a noun and so on. These authors then enable the possibility of \emph{uninflectable} entries inside the group of otherwise inflecting parts of speech in Japanese, an analysis which opens up room for discussion.

\subsection{Possible \emph{na}-adjective \emph{rentaishi}\is{adjective!Japanese|(}} \label{12sec5.1}

This discussion will be lead here, as it is most eminent, with regard to \emph{rentaishi} which are arguably derived from adjectives and are thus argued by the above-mentioned authors to belong to this group. First and foremost, this concerns the three lexemes \emph{ōki-na} `big',
\emph{chiisa-na} `small' and \emph{okashi-na} `strange' (or potentially \emph{ōkina, chiisana} and \emph{okashina}), which all end in \emph{-na}, thus look like \emph{na}-adjectives. However, they do not show inflection, in other words lack predicative use with the copula \emph{da} and adverbial use with the element \emph{-ni}. This is illustrated for \emph{ōki-na} `big' below.\footnote{To illustrate potential inflection, I will transliterate the ending \emph{-na} with a hyphen.}

\begin{exe}
	\ex \label{12ex40}
	\begin{xlist}
		\ex 
		\gll ōki-na ie \\
		big-\textsc{na} house \\
		\trans ‘big house’
		
		\ex[*] {
			\gll Kono ie-wa ōki-da. \\
			this house-\textsc{top} big-\textsc{cop} \\
			\trans Intend.: ‘This house is big.’}
		
		\ex[*] {
			\gll Ie-wa ōki-ni nat-ta. \\
			house-\textsc{top} big-\textsc{adv} become-\textsc{pst} \\
			\trans Intend.: ‘The house became big.’}
	\end{xlist}
\end{exe}

% \largerpage[1]

However, there exists \emph{i}-adjective variants of these lexemes that are fully productive and do not adhere to these restrictions. Crucially, the meaning of the attributive modifier remains the same. This is shown below.

\begin{exe}
	\ex \label{12ex41}
	\begin{xlist}
		\ex 
		\gll ōki-i ie \\
		big-\textsc{i} house \\
		\trans ‘big house’
		
		\ex 
		\gll Kono ie-wa ōki-i. \\
		this house-\textsc{top} big-\textsc{prs} \\
		\trans ‘This house is big.’
		
		\ex 
		\gll Ie-wa ōki-ku nat-ta. \\
		house-\textsc{top} big-\textsc{adv} become-\textsc{pst} \\
		\trans ‘The house became big.’
	\end{xlist}
\end{exe}

There are three possible ways to analyze these lexemes in terms of word class status. We can either classify them as \emph{rentaishi}, thus as an uninflecting category, presupposing that they somehow have become fossilized. This explanation seems reasonable, although it is hard to determine at what point of the language this fossilization occurred. The first result for \emph{ōki-na} in the Corpus of Historical Japanese (CHJ\footnote{\url{https://clrd.ninjal.ac.jp/chj/overview-en.html}}) stems from 1642, thus from Middle Japanese. Furthermore, it is unclear if they are derived from their \emph{i}-adjective counterparts or potentially derived from forms such as \emph{ōki-naru} in which a lexical adjective is combined with the historical copula \emph{nari}. The first result for the form \emph{ōki-naru} is attested on the CHJ in 1001.

Before we discuss their incorporation into the \emph{na}-adjective category, another prominent option needs to be discussed, namely the analysis as \emph{i}-adjectives with a secondary attributive form \emph{-na}, which seems to be the preferred option in the research literature (\citealt{suzuki1972, rickmeyer1995, muraki2012, lehmannnishina2015} among others). One piece of evidence presented by authors such as \citet{suzuki1972} and \citet{lehmannnishina2015} in favor of this analysis is that there is no semantic difference between the respective forms. Therefore, there should be no need for positing two separate lexical entries in two different word classes.\footnote{\citet{lehmannnishina2015} also count lexemes such as \emph{atataka-na} `warm' as alternative forms of existing \emph{i}-adjectives. However, such lexemes are not syntactically restricted and should thus be treated as true \emph{na}-adjectives. Japanese then possibly simply has two separate lexical entries across two word classes in such cases.} Although the existence of uninflectable lexemes, and alternative inflectional forms for that matter, is not foreseen according to the more traditional analyses of Japanese word classes depicted above, this option is justified, considering the semantic closeness of the two forms. In fact, whether we allow the existence of uninflectable lexemes in Japanese is close to an ideological debate and especially with regard to these three lexemes, nothing hinges on this distinction (see also \citealt[Fn. 15]{oshimahayashi2021-nihongo}).

This brings us to the third alternative, namely to classify these lexemes as \emph{uninflectable} or \emph{defective na}-adjectives. This approach will be discussed below.

\subsection{\emph{Rentaishi} as defective adjectives\is{defectiveness!syntactical|(}} \label{12sec5.2}

The most contemporary understanding of \emph{rentaishi} is equal to `adjectives without predicative use'. This view emerged in the past 20 years and is found most prominently in the works of Muraki Shinjirō (for example \citeyear{muraki1996}, \citeyear{muraki2009}, \citeyear{muraki2012}, \citeyear{muraki2015}, see also \citealt{oshimaakitasano2019} and \citealt{oshimahayashi2021-nihongo, oshimahayashi2021-redefining}). Muraki, in his role as important adversary of the school grammar, uses \emph{rentaishi} synonymously to `mono-attributive adjectives' (\emph{kitei-keiyōshi}), that is adjectives which can only be used attributively. He argues in favor of three morphologically diverse adjective groups, ending in \emph{-i}, \emph{-na} and \emph{-no} respectively, and emphasizes that all adjectives should in principle be able to occur attributively, predicatively and adverbially \citet[36--38]{muraki2012}. If an adjectival lexeme lacks one of these uses it is described as ``defective'', or literally ``incomplete'' (\emph{fukanzen}, \citealt[37]{muraki2012}; \citealt[26]{muraki2015}). One important group of incomplete adjectives are those that are mono-attributive, and those mono-attributive adjectives are what he calls \emph{rentaishi}. Under this view, the three lexemes discussed in the previous section are \emph{na-rentaishi}.

Importantly, however, most mono-attributive (non-predicative) modifiers in Japanese can be found among the group of modifiers co-occurring with the element \emph{-no}, named \emph{no}-modifiers\is{Japanese \emph{no}-modifier} in \citet{kohlich2023}, but analyzed as \emph{no}-adjectives\is{Japanese \emph{no}-modifier} in \citet{muraki2012}, as evidenced in recent quantitative surveys on Japanese modifiers performed on the Balanced Corpus of Contemporary Written Japanese (BCCWJ\footnote{\url{https://clrd.ninjal.ac.jp/bccwj/en/}}) \citep{abekitazawaliu2022, kohlich2023}. In the corpus analysis of \citet{abekitazawaliu2022}, all modifiers given in \citet{muraki2012} are analyzed for the variables \emph{attributive} and \emph{predicative} \emph{use} (among others). The authors determine 12 lexemes that completely lack predicative use, which they classify accordingly as \emph{rentaishi}.\footnote{More lexemes, around 300, are classified as non-predicative in \citet{kohlich2023} which is due to the facts that more lexemes are analyzed in general and that also lexemes between one and five occurrences, or a ratio of 1:25 towards attributive use are classified as non-predicative modifiers.}

\largerpage[0.5]

\citet{oshimaakitasano2019} and \citet{oshimahayashi2021-nihongo, oshimahayashi2021-redefining} go in the same direction, although for them \emph{rentaishi} are not necessarily adjectives. If a lexeme appears attributively with \emph{-no} but cannot be used predicatively they call it `\emph{no}-Type Adnominal', or `\emph{no-gata rentaishi}', lit. `\emph{no}-type \emph{rentaishi}'. They give the following definition.

\begin{exe}
	\ex \label{12ex42}
	Those lexemes that form a noun-modifying clause being accompanied by the attributive copula form \emph{no}, but cannot be accompanied by other copula forms such as \emph{da} and \emph{de}. 
	
	\hfill \citep[520]{oshimahayashi2021-redefining}
\end{exe}

The authors analyze \emph{rentaishi} as compositional, as a combination of a lexeme and the element \emph{-no}, which they define as the copula.\footnote{One needs to abstract away from the fact that the authors understand \emph{no} as the copula when they talk about non-predicative lexemes. This analysis is of course not tenable as non-predicative elements should not be able to co-occur with the copula.}

To summarize, for the authors mentioned above, the term `\emph{rentaishi}' does not refer to an uninflecting word class\is{uninflecting!word class|(}, but is perceived as a feature, namely that of being non-predicative or [-predicative]. Under this view, the lexemes in question are \emph{defective} in the sense of \citet{spencer2020} as they simply do not occur in those syntactic environments they are expected to appear in if they are treated as adjectives. They are thus on a par with \ili{English} \emph{former}\is{adjective!non-intersective} or \ili{Russian} \emph{kni\v{z}nyj} `pertaining to books', as mentioned in Section \ref{12sec2}.

This potentially opens a new avenue for future research, as the connection of \isi{uninflectedness} and uninflectability respectively to restricted uses of adjectives to the attributive or predicative domain has, to my knowledge, not been extensively discussed in the research literature.\footnote{I express my gratitude to an anonymous reviewer who pointed out this issue to me and lead me to rethink my, formerly rigid, opinion on this matter.} Another cross-linguistic case where such a discussion could be relevant are mono-predicative adjectives in \ili{German}. While \ili{German} adjectives, in general, only show inflection in attributive, never in predicative, position there are some adjectival lexemes such as \textit{barfuß} ‘barefoot’ or \textit{pleite} ‘bankrupt, broke’, which are strongly biased toward predicative use and show highly restricted or marginal attributive distribution \citep{trost2006}. Treating the possibility of inflection as a criterion for adjectivehood, those lexemes are not adjectives, but could be perceived as some reversed counterpart to Japanese \emph{rentaishi}. On the other hand, their uninflectability is connected to their syntactic behavior. Such lexemes should be studied with regard to the debate at hand.

While the view of the term \emph{rentaishi} as mono-attributive adjectives in Japanese is appealing, I would like to conclude this section with two remaining problems this view entails. First, we are in danger of losing the historical dimension of the notion \emph{rentaishi}. Lexemes classified as \emph{rentaishi} in earlier approaches can be found at very early stages of Japanese. For example, \emph{reino} `ordinary, aforementioned' can be traced back at least to
the Heian-period (794--1185). The first result on the CHJ dates to the year 900. This is different for other non-predicative \emph{no}-modifiers\is{Japanese \emph{no}-modifier}, such \emph{as bōgai-no} `unexpected' which is classified as a \emph{rentaishi} in \citet{abekitazawaliu2022} but for which the first entry on the CHJ dates back only to 1887. From a historical point of view, namely following an analysis of the word class \emph{rentaishi} as fossilized derivatives, it might therefore be difficult to argue for membership of lexemes such as \emph{bōgai-no} to the word class \emph{rentaishi}. This would thus be one problem one has to live with under this more modern approach.

A bigger problem specific to the work of \citeauthor{muraki2012} (\citeyear{muraki2012} among others) regards the fact that this author does not recognize verbally derived mono-attributive lexemes such as \emph{aru} `(a) certain' as \emph{rentaishi}, as he only discusses adjectival lexemes. However, I do not see a reason why only adjectives should be treated as having a possible restriction to attributive or predicative position in the first place.

I propose to maintain the historical dimension of \emph{rentaishi} as an uninflecting word class\is{uninflecting!word class|)} whose members only serve as adnominal modifiers. On the other hand, I think it is appropriate to extend the notion \emph{rentaishi} to refer to lexemes from other word classes that lack predicative use even if they do not belong to this word class from a historical point of view. One consequence of this is that the group of \emph{rentaishi} becomes an open word class, although linguists in the Meiji-period probably did not perceive the word class \emph{rentaishi} as open. Another crucial consequence of this is that the morphological heterogeneity of the word class \emph{rentaishi} falls out nicely as different morphological subgroups are not only allowed, but predicted\is{uninflectable!element|)}\is{defectiveness!syntactical|)}\is{adjective!Japanese|)}.

\section{Summary} \label{12sec6}

This paper looked at Japanese with regard to the debate on \emph{\isi{uninflectedness}}. Although the notion of \emph{inflection}, in Japanese \emph{katsuyō}\is{katsuyō}, is slightly different from what we might expect from our knowledge of Indo-European languages, this notion has been used as the basis of word class classification in Japan roughly since the 17\textsuperscript{th} century. Verbs and adjectives are classified as categorically inflecting, contra nouns as categorically uninflecting parts of speech.

It was shown that from the beginning of the 20\textsuperscript{th} century, linguists in Japan became aware of the differences between their adjectives and those of Western languages, mostly \ili{English} and \ili{Dutch}. They realized that adjectives in those languages cannot inflect independently, but only in accordance with their head noun. Having realized this difference, they defined the Japanese adjective, the \emph{keiyōshi}, as more powerful.

At the same time, they realized that there are some lexemes in Japanese whose sole purpose is attributive modification. In order to fill the slot of dependent and defective attributive modifiers, they literally created the word class of \emph{rentaishi} as the Japanese counterparts to the Western adjective and sorted all those attributive modifiers into this group. This was famously done in the grammar of \citeauthor{matsushita1976} Daizaburō (\citeyear{matsushita1976} [1928]) and the school grammar\is{Japanese school grammar} of \citeauthor{hashimoto1948} Shinkichi (\citeyear{hashimoto1948} [1934]) allowing for the emergence of \emph{rentaishi} as an independent word class.

In a nutshell, \emph{rentaishi} present a case of an \emph{uninflecting} word class established on historical grounds. It was exactly the characteristic \emph{\isi{uninflectedness}} which served as the motivation for word class status of relevant lexemes. From a modern perspective, this word class, as it was perceived at the beginning of the 20\textsuperscript{th} century, arguably poses an example of a closed category of only a few lexemes, defined via the characteristic \emph{\isi{uninflectedness}}.

Towards the end of this paper, it was outlined that starting around the middle of the 20\textsuperscript{th} century, several authors performed a re-analysis of \emph{rentaishi} and, instead, sorted relevant lexemes into the word classes they morphologically belong to. Nowadays, the notion \emph{rentaishi} is mostly treated as a syntactic feature, synonymous to `mono-attributive' or [-predicative]. I argued that, if we are able to maintain the historical definition of \emph{rentaishi}, this view can open a new avenue of research.

For future research, it would thus be interesting to tie the debate on \isi{uninflectedness} to the discussion of adjectives without attributive use and to extend this approach to those without predicative use. Furthermore, it is desirable to test the characteristics inflecting/uninflecting or inflectable/uninflectable against the notion of word classes in other languages and to compare them, alongside any historical relevant developments, to the Japanese case\il{Japanese|)}\is{rentaishi|)}.
 
 
\section*{Abbreviations}

\begin{tabularx}{.45\textwidth}{lQ}
\textsc{acc} & accusative \\
\textsc{com} & comitative \\
\textsc{comp} & complementizer \\
\textsc{cop} & dummy copula \\
\textsc{imp} & imperative \\
\textsc{n} & neuter \\
\textsc{neg} & negative \\
\textsc{nom} & nominative \\
\textsc{pred} & predictive copula \\
\textsc{prep} & prepositional case
\end{tabularx}
\begin{tabularx}{.45\textwidth}{lQ}
\textsc{prs} & non-past \\
\textsc{pst} & past \\
\textsc{pst.ptcp} & past participle \\
\textsc{ra} & relational adjective \\
\textsc{rc} & relative clause \\
\textsc{sg} & singular \\
\textsc{subj} & subjunctive \\
\textsc{top} & topic \\
\textsc{vol} & volitional
\end{tabularx}

%\section*{Acknowledgements}

%\section*{Contributions}
%John Doe contributed to conceptualization, methodology, and validation.
%Jane Doe contributed to writing of the original draft, review, and editing.


{\sloppy\printbibliography[heading=subbibliography,notkeyword=this]}
\end{document}
